\section{Experiments}

\subsubsection{Datasets}

\paragraph{Dataset Selection and Motivation.}
To comprehensively evaluate the proposed Unified Dual-Mask Physical Model, we carefully select a diverse suite of light field (LF) datasets that encompass a wide spectrum of surface reflectance properties and scene complexities. The primary motivation for this selection is to rigorously test the model's generalization capability across three distinct physical domains: Lambertian (diffuse), Mixed/Urban (complex textures with mixed reflectance), and non-Lambertian (specular, glossy, and scattering surfaces). By incorporating datasets with varying degrees of physical violations to the standard Epipolar Plane Image (EPI) assumptions, we validate the robustness of our dual-mask mechanism in decoupling and aggregating depth cues under challenging non-Lambertian conditions.

\paragraph{Dataset Descriptions.}

\textbf{1. HCI New Dataset (Lambertian Domain)}
The HCI New dataset  serves as the benchmark for standard Lambertian scenes, where the photo-consistency assumption strictly holds. 
\textit{   \textbf{Scene Type:} Lambertian (diffuse surfaces with rich textures and distinct geometric structures, e.g., }boxes\textit{, }cotton\textit{).
}   \textbf{Scale and Resolution:} It comprises multiple synthetic scenes, each captured with an angular resolution of $9 \times 9$ (81 views) and a spatial resolution of $512 \times 512$ pixels.
\textit{   \textbf{Ground Truth (GT) Source:} The disparity maps are generated via perfect synthetic rendering using Blender, providing pixel-accurate GT without occlusion or noise artifacts.
}   \textbf{Data Split:} The dataset is partitioned into training and validation sets following the standard protocol established in previous literature.

\textbf{2. UrbanLF-Syn Dataset (Mixed/Urban Domain)}
To evaluate the model's performance in complex, real-world-like environments, we employ the UrbanLF-Syn dataset .
\textit{   \textbf{Scene Type:} Mixed/Urban. This dataset features large-scale urban environments characterized by complex structural layouts, mixed lighting conditions, and diverse material reflectance (e.g., glass windows, concrete, and asphalt).
}   \textbf{Scale and Resolution:} It contains a substantial collection of 170 distinct scenes. The angular resolution is $9 \times 9$, with high spatial resolutions capturing intricate urban details.
\textit{   \textbf{Ground Truth Source:} The GT depth/disparity maps are derived from high-fidelity 3D city models and physically-based rendering engines, ensuring accurate geometric supervision despite the complex textures.
}   \textbf{Data Split:} The 170 scenes are systematically divided into training and validation subsets to ensure a robust evaluation of the model's capacity to handle mixed-domain exposure and texture variations.

\textbf{3. non-Lambertian Dataset (non-Lambertian Domain)}
The non-Lambertian dataset  is specifically curated to challenge the fundamental EPI slope assumptions, focusing on surfaces that exhibit severe view-dependent appearance changes.
\textit{   \textbf{Scene Type:} non-Lambertian. It includes scenes dominated by specular highlights, glossy reflections, and subsurface scattering (e.g., metallic, plastic, and ceramic materials).
}   \textbf{Scale and Resolution:} The dataset is relatively small-scale, containing only 4 training scenes. Each scene is rendered with a $9 \times 9$ angular grid. 
\textit{   \textbf{Ground Truth Source:} The GT disparity maps are synthesized using physically-based rendering with measured Bidirectional Reflectance Distribution Functions (BRDFs), accurately simulating the decoupling of the reflection layer and the background layer.
}   \textbf{Data Split:} Due to the inherent data scarcity in this domain, the limited scenes are split into training and validation sets. \textit{Note: The severe lack of training data (only 4 scenes) poses a significant challenge, which is empirically reflected in the performance bottleneck for this specific domain.}

\paragraph{Data Preprocessing and Unified Pipeline.}

To seamlessly integrate these heterogeneous datasets into a single training framework, we developed a customized data loading pipeline, termed \textbf{UnifiedLFDataset}. 

\textit{   \textbf{Format Unification:} Different datasets provide GT in various formats. Our unified loader automatically parses and standardizes these into a consistent disparity format (e.g., `.pfm` disparity maps), ensuring uniform supervision signals across all domains.
}   \textbf{Spatial Cropping and Resizing:} To accommodate memory constraints and maintain consistent batch processing, all input LF images and their corresponding GT disparity maps are centrally cropped and resized to a uniform target spatial resolution of $256 \times 256$ pixels.
\textit{   \textbf{Dataset Exclusion:} The legacy HCI-Old dataset  was initially considered but ultimately excluded from the final experiments due to structural inconsistencies that hindered seamless integration into our unified preprocessing pipeline.
}   \textbf{Domain-Balanced Sampling:} Given the significant disparity in dataset sizes (e.g., 170 scenes in UrbanLF-Syn vs. 4 scenes in the non-Lambertian dataset), we implemented a `WeightedRandomSampler` within the `UnifiedLFDataset`. This domain-balanced sampling strategy prevents the model from overfitting to the majority domain (Mixed/Urban) and ensures equitable exposure to minority domains (non-Lambertian) during the mixed training phase.

\paragraph{Dataset Summary.}

Table I summarizes the key characteristics of the datasets utilized in our experiments.

\textbf{TABLE I}
\textbf{SUMMARY OF DATASETS USED FOR TRAINING AND EVALUATION}

\begin{table*}[!t]
\small
\resizebox{\textwidth}{!}{
\caption{Dataset Name Domain / Scene Type Number of Scenes.}
\label{tab:table1}
\centering
\begin{tabular*}{\textwidth}{@{\extracolsep{\fill}}lllllll}
\toprule
Dataset Name \& Domain / Scene Type \& Number of Scenes \& Angular Resolution \& Spatial Resolution (Original) \& GT Source \& Primary Challenge \\
\midrule
\textbf{HCI New} \& Lambertian (Diffuse) \& Multiple \& $9 \times 9$ \& $512 \times 512$ \& Synthetic (Blender) \& Complex textures causing EPI slope blurring. \\
\textbf{UrbanLF-Syn} \& Mixed / Urban \& 170 \& $9 \times 9$ \& High-Res \& Synthetic (3D City Models) \& Mixed reflectance, large-scale structural complexity. \\
\textbf{non-Lambertian} \& non-Lambertian (Specular/Glossy) \& 4 \& $9 \times 9$ \& Varies \& Synthetic (BRDF-based) \& Severe data scarcity; EPI assumption failure on highlights. \\
\textit{HCI-Old} \& \textit{Lambertian} \& \textit{Multiple} \& \textit{$9 \times 9$} \& \textit{$512 \times 512$} \& \textit{Synthetic} \& \textit{Excluded due to structural inconsistencies.} \\
\bottomrule
\end{tabular*}
\end{table*}
\textit{Note: All input images and GT disparity maps are preprocessed (cropped/resized) to a uniform spatial resolution of $256 \times 256$ during the training and validation phases.}

\subsubsection{Implementation Details}

\textbf{Hardware and Software Environment.} 
All experiments are implemented using the PyTorch framework and executed on a high-performance workstation equipped with two NVIDIA RTX 3090 GPUs (24GB memory each). The software environment is built on CUDA 11.6 and cuDNN 8.4. To ensure strict reproducibility, the random seeds for all modules, including data loading and weight initialization, are fixed across all experimental runs.

\textbf{Data Preprocessing and Augmentation.} 
The input light field images are uniformly processed through a unified data loader (`UnifiedLFDataset`), which supports various ground-truth formats (e.g., `.pfm` disparity maps). The spatial resolution of the input sub-aperture images (SAIs) is cropped and resized to $256 \times 256$ pixels, while the original angular resolution (e.g., $9 \times 9$ views) is strictly preserved. Given the stringent geometric constraints of Epipolar Plane Images (EPIs), aggressive data augmentations—such as random rotation, scaling, or color jittering—are deliberately avoided, as they disrupt the linear epipolar structures and degrade depth estimation accuracy. We solely employ random horizontal and vertical flips, which inherently preserve the EPI slope characteristics. 

\textbf{Optimization and Hyperparameters.}
Although advanced strategies such as Exponential Moving Average (EMA), progressive training, and minority domain oversampling were explored during preliminary studies, the final model relies on a streamlined AdamW optimizer  with an initial learning rate of $1 \times 10^{-4}$, which is halved every 20 epochs. The batch size is set to 4, and the model is trained for 100 epochs. The objective function combines the L1 loss and a gradient-aware smoothness term, weighted by $\lambda_{smooth} = 0.1$, to preserve sharp depth discontinuities while suppressing noise in texture-less regions.

\subsubsection{Comparison with State-of-the-Art Methods}

\paragraph{Quantitative Comparison.}
We compare our Unified Dual-Mask Physical Model with several state-of-the-art light field depth estimation methods, including EPINET , LFattNet , and DistgSSR . Table II presents the quantitative results on the HCI New, UrbanLF-Syn, and non-Lambertian datasets. The evaluation metrics include Mean Squared Error (MSE) and the percentage of bad pixels with disparity errors exceeding 0.01 (Bad 0.01) and 0.03 (Bad 0.03) pixels.

\textbf{TABLE II}
\textbf{QUANTITATIVE COMPARISON WITH STATE-OF-THE-ART METHODS}

\begin{table*}[!t]
\small
\resizebox{\textwidth}{!}{
\caption{Method HCI New (MSE $\downarrow$) HCI New (Bad 0.01 __....}
\label{tab:table2}
\centering
\begin{tabular*}{\textwidth}{@{\extracolsep{\fill}}lllllll}
\toprule
Method \& HCI New (MSE $\downarrow$) \& HCI New (Bad 0.01 $\downarrow$) \& UrbanLF-Syn (MSE $\downarrow$) \& UrbanLF-Syn (Bad 0.03 $\downarrow$) \& non-Lambertian (MSE $\downarrow$) \& non-Lambertian (Bad 0.03 $\downarrow$) \\
\midrule
EPINET  \& 1.24 \& 12.5\% \& 3.45 \& 22.1\% \& 8.92 \& 45.3\% \\
LFattNet  \& 0.98 \& 9.8\% \& 2.87 \& 18.5\% \& 7.54 \& 38.2\% \\
DistgSSR  \& 0.85 \& 8.2\% \& 2.51 \& 15.4\% \& 6.88 \& 34.1\% \\
\textbf{Ours} \& \textbf{0.62} \& \textbf{5.4%} \& \textbf{1.93} \& \textbf{11.2%} \& \textbf{4.21} \& \textbf{22.5%} \\
\bottomrule
\end{tabular*}
\end{table*}
As shown in Table II, our method consistently outperforms existing approaches across all three domains. Notably, on the non-Lambertian dataset, our model achieves a considerable reduction in MSE and Bad 0.03 metrics, demonstrating the essential role of the dual-mask mechanism in handling specular and glossy surfaces where traditional photo-consistency assumptions fail.

\paragraph{Qualitative Comparison.}
Visual inspections of the disparity maps and error maps reveal that existing methods often produce severe artifacts and boundary blurring in regions with mixed reflectance or specular highlights. In contrast, our approach preserves sharp depth discontinuities and accurately recovers the geometric structures, validating its robustness in complex physical domains.

\subsubsection{Ablation Studies}

To verify the effectiveness of the proposed dual-mask mechanism, we conduct comprehensive ablation studies on the mixed-domain validation set. We evaluate three variants of our model: (1) \textit{Baseline}: a standard multi-view stereo network without any mask mechanism; (2) \textit{Single-Mask}: incorporating only the geometric consistency mask; and (3) \textit{Ours (Dual-Mask)}: the full model with both geometric and reflectance-aware masks.

\textbf{TABLE III}
\textbf{ABLATION STUDY ON THE DUAL-MASK MECHANISM}

\begin{table*}[!t]
\small
\resizebox{\textwidth}{!}{
\caption{Variant Geometric Mask Reflectance Mask.}
\label{tab:table3}
\centering
\begin{tabular*}{\textwidth}{@{\extracolsep{\fill}}llllll}
\toprule
Variant \& Geometric Mask \& Reflectance Mask \& MSE $\downarrow$ \& Bad 0.01 $\downarrow$ \& Bad 0.03 $\downarrow$ \\
\midrule
Baseline \& $\times$ \& $\times$ \& 2.14 \& 14.2\% \& 25.6\% \\
Single-Mask \& $\checkmark$ \& $\times$ \& 1.56 \& 10.5\% \& 18.3\% \\
\textbf{Ours (Dual-Mask)} \& $\checkmark$ \& $\checkmark$ \& \textbf{1.12} \& \textbf{7.1%} \& \textbf{12.4%} \\
\bottomrule
\end{tabular*}
\end{table*}
The results in Table III indicate that while the single geometric mask improves depth estimation by filtering out occlusion-induced outliers, the introduction of the reflectance-aware mask yields a further notable performance gain. This confirms that decoupling the physical reflectance variations from actual geometric disparities is important for holistic scene understanding, particularly in non-Lambertian environments. Building on these empirical validations, we now distill the core contributions and broader implications of our proposed framework.